% 来源：sources/2016-2025-期中-甲-历年真题汇编.pdf 第 5 页
\begin{problem}{15 分}{8.5cm}
\[
D=\begin{vmatrix}
a_1+x&a_2&\cdots&a_n\\
a_1&a_2+x&\cdots&a_n\\
\vdots&\vdots&&\vdots\\
a_1&a_2&\cdots&a_n+x
\end{vmatrix}.
\]
\end{problem}

\begin{problem}{15 分}{8.5cm}
设 $\sigma=(i_1i_2\cdots i_n)$ 是一个 $n$ 阶排列，$A=[a_{ij}]$ 是一个 $n$ 阶方阵，并且 $A$ 中元素满足：对于每个固定 $r$，当 $j=i_r$ 时，$a_{rj}=1$，否则 $a_{rj}=0$，求 $|A|$。
\end{problem}

\begin{problem}{15 分}{9cm}
问：当 $\lambda$ 取什么值时，方程组无解？唯一解？无穷多解？有解时求解。
\[
\begin{cases}
\lambda x_1+x_2+x_3=1,\\
x_1+\lambda x_2+x_3=1,\\
x_1+x_2+\lambda x_3=1.
\end{cases}
\]
\end{problem}

\begin{problem}{15 分}{8.5cm}
设 $A=\begin{pmatrix}1&1&1\\1&1&1\\1&1&1\end{pmatrix}$，求矩阵方程：$AX=A+2X$。
\end{problem}

\begin{problem}{15 分}{8.5cm}
设 $A$ 为 $4$ 阶反对称矩阵，$B=\begin{pmatrix}0&0&0&0\\0&0&0&0\\0&0&1&0\\0&0&0&2\end{pmatrix}$，证明：$E+AB$ 可逆。
\end{problem}

\begin{problem}{10 分}{7cm}
设 $A$ 为 $n$ 阶实对称矩阵 $(n>1)$，$|A|=0$，证明：$A_{ij}A_{ji}=(A_{ij})^{2}$，$(i,j=1,2,\cdots,n)$。
\end{problem}

\begin{problem}{7 分}{6.5cm}
$A,B\in P^{n\times n}$，$M=\begin{pmatrix}A&B\\B&A\end{pmatrix}$，证明：$R(M)\ge R(A+B)+R(A-B)$。
\end{problem}

\begin{problem}{8 分}{6.5cm}
设 $A,B\in P^{n\times n}$，满足 $B=E+AB$，证明：$AB=BA$。
\end{problem}
